% !TeX program = xelatex
% UTF-8；用 XeLaTeX 编译两次。仅为模板示例，未转换原讲义。
\documentclass[aspectratio=169,11pt]{beamer}
\usepackage[UTF8,fontset=fandol]{ctex}
\usepackage{amsmath,amssymb,bm,booktabs,array}
\usepackage{tikz,pgfplots}
\pgfplotsset{compat=1.18}
\usetikzlibrary{calc}

% ===== 修改这里即可更换整套幻灯片的颜色 =====
\definecolor{Navy}{HTML}{17354D}
\definecolor{Teal}{HTML}{167D8D}
\definecolor{Ink}{HTML}{263746}
\definecolor{Muted}{HTML}{71808D}
\definecolor{Wash}{HTML}{F0F5F7}
\definecolor{Rule}{HTML}{DCE5EA}
\definecolor{Gold}{HTML}{B88B48}

\usetheme{default}
\usefonttheme{professionalfonts}
\setbeamersize{text margin left=10mm,text margin right=10mm}
\setbeamercolor{normal text}{fg=Ink,bg=white}
\setbeamercolor{structure}{fg=Teal}
\setbeamercolor{frametitle}{fg=Navy}
\setbeamercolor{title}{fg=Navy}
\setbeamercolor{subtitle}{fg=Muted}
\setbeamercolor{block title}{fg=Navy,bg=Wash}
\setbeamercolor{block body}{fg=Ink,bg=Wash}
\setbeamercolor{alerted text}{fg=Teal}
\setbeamerfont{title}{size=\LARGE,series=\bfseries}
\setbeamerfont{subtitle}{size=\normalsize}
\setbeamerfont{frametitle}{size=\Large,series=\bfseries}
\setbeamerfont{block title}{size=\normalsize,series=\bfseries}
\setbeamertemplate{navigation symbols}{}
\setbeamertemplate{headline}{}
\setbeamertemplate{items}[circle]
\setbeamertemplate{itemize subitem}{\tiny$\circ$}
\setlength{\parskip}{0.35em}

\setbeamertemplate{frametitle}{%
  \vspace{4mm}\insertframetitle\par
  \vspace{1.5mm}{\color{Rule}\hrule height 0.5pt}\vspace{2mm}}
\setbeamertemplate{footline}{%
  \leavevmode\hbox{%
  \begin{beamercolorbox}[wd=.80\paperwidth,ht=4mm,dp=2.5mm,leftskip=10mm]{normal text}
    \usebeamerfont{footline}\color{Muted}\scriptsize\insertshorttitle
  \end{beamercolorbox}%
  \begin{beamercolorbox}[wd=.20\paperwidth,ht=4mm,dp=2.5mm,rightskip=10mm plus 1fil]{normal text}
    \hfill\color{Muted}\scriptsize\insertframenumber\,/\,\inserttotalframenumber
  \end{beamercolorbox}}}

\newcommand{\smallnote}[1]{\par{\scriptsize\color{Muted}#1\par}}
\newcommand{\labeltext}[1]{{\color{Teal}\bfseries #1}\par}
\newcommand{\spaceditem}{\setlength{\itemsep}{0.8em}}

% ===== 报告信息：将下列字段替换为自己的信息 =====
\title[空间数据与机器学习 · 模板示例]{空间数据与机器学习}
\subtitle{学术汇报模板示例\\[2mm]A Beamer Template for Research Talks}
\author{汇报人姓名}
\institute{学校 / 学院 / 课题组}
\date{2026 年 10 月}
\hypersetup{pdftitle={Beamer 学术汇报模板示例},pdfauthor={模板示例}}

\begin{document}

% 01 / 封面
\begin{frame}[plain]
  \vspace{7mm}
  {\color{Teal}\small 学术汇报 \quad RESEARCH SEMINAR\par}
  \vspace{9mm}
  {\usebeamerfont{title}\color{Navy}\inserttitle\par}
  \vspace{4mm}
  {\usebeamerfont{subtitle}\color{Muted}\insertsubtitle\par}
  \vspace{7mm}
  {\color{Teal}\rule{20mm}{1.2pt}\par}
  \vspace{4mm}
  {\normalsize\insertauthor\par}
  {\small\color{Muted}\insertinstitute\par}
  \vspace{1mm}
  {\small\color{Muted}\insertdate\par}
  \vfill
  \smallnote{本文件演示幻灯片版式。公式为教学示例，图表数值为合成数据。}
\end{frame}

% 02 / 报告结构：条目简短，保留讲述空间
\begin{frame}{报告结构}
  \vspace{2mm}
  \begin{tabular}{@{}>{\color{Teal}\Large\bfseries}p{.09\textwidth}
                    >{\raggedright\arraybackslash}p{.81\textwidth}@{}}
    01 & \textbf{研究问题}\\[-1mm]
       & {\small\color{Muted}数据包含什么，希望解释什么}\\[2.5mm]
    02 & \textbf{模型与方法}\\[-1mm]
       & {\small\color{Muted}变量定义、目标函数与建模假设}\\[2.5mm]
    03 & \textbf{实验与分析}\\[-1mm]
       & {\small\color{Muted}结果图、对比表与证据边界}\\[2.5mm]
    04 & \textbf{结论与讨论}\\[-1mm]
       & {\small\color{Muted}主要观点、局限与后续问题}
  \end{tabular}
\end{frame}

% 03 / 双栏背景页
\section{研究问题}
\begin{frame}{问题设定：表达信息与空间信息}
  \begin{columns}[T,onlytextwidth]
    \begin{column}{.47\textwidth}
      \labeltext{表达矩阵}
      \[
        X\in\mathbb{R}^{n\times p}
      \]
      \begin{itemize}\spaceditem
        \item $n$ 个观测位置
        \item $p$ 个测量特征
        \item 每行对应一个位置
      \end{itemize}
    \end{column}
    \begin{column}{.47\textwidth}
      \labeltext{空间坐标}
      \[
        S\in\mathbb{R}^{n\times 2}
      \]
      \begin{itemize}\spaceditem
        \item 二维平面中的观测位置
        \item 邻接矩阵 $A\in\mathbb{R}^{n\times n}$
        \item $A$ 描述空间邻接关系
      \end{itemize}
    \end{column}
  \end{columns}
  \vspace{0mm}
  \begin{block}{研究问题}
    如何在低维表示中同时考虑表达信息与空间邻接关系？
  \end{block}
  \vspace{1mm}
  \smallnote{教学设定：变量与维度用于展示双栏、符号和重点结论的排版。}
\end{frame}

% 04 / 核心公式页
\section{模型与方法}
\begin{frame}{方法示例：带空间正则项的低秩表示}
  \vspace{1mm}
  \[
    \min_{\substack{Z,W\\W^{\mathsf T}W=I_d}}\quad
    \underbrace{\lVert X-ZW^{\mathsf T}\rVert_F^2}_{\text{表达重构}}
    \;+\;
    \lambda\underbrace{\operatorname{tr}(Z^{\mathsf T}LZ)}_{\text{空间正则}}
  \]
  \vspace{5mm}
  \begin{columns}[T,onlytextwidth]
    \begin{column}{.48\textwidth}
      \labeltext{变量与约束}
      \begin{itemize}\spaceditem
        \item $Z\in\mathbb{R}^{n\times d}$：低维表示
        \item $W\in\mathbb{R}^{p\times d}$：特征载荷
        \item $d\ll \min(n,p)$，$\lambda\geq 0$
      \end{itemize}
    \end{column}
    \begin{column}{.48\textwidth}
      \labeltext{空间项的含义}
      \[
        L=D-A,\qquad D_{ii}=\sum_j A_{ij}
      \]
      \small 对称非负邻接矩阵下，空间正则项惩罚相邻位置的表示差异。
    \end{column}
  \end{columns}
  \vfill
  \smallnote{该目标函数仅为排版与建模示例，不对应讲义中的特定方法或实验。}
\end{frame}

% 05 / 图表：原生 PGFPlots，数据直接编辑
\section{实验与分析}
\begin{frame}{结果图示例：目标函数的变化}
  \begin{columns}[c,onlytextwidth]
    \begin{column}{.68\textwidth}
      \begin{tikzpicture}
      \begin{axis}[
        width=.96\linewidth,height=49mm,
        xmin=0,xmax=50,ymin=0,ymax=1.05,
        xlabel={迭代次数},ylabel={归一化目标函数值},
        label style={font=\small},tick label style={font=\scriptsize},
        xtick={0,10,20,30,40,50},ytick={0,.25,.5,.75,1},
        axis lines=left,axis line style={Rule},
        tick style={Rule},grid=major,grid style={Rule,dashed},
        legend style={at={(.97,.97)},anchor=north east,
          draw=none,fill=white,font=\scriptsize},
        clip=true]
        \addplot[Navy,thick,mark=*,mark size=1.6pt]
          coordinates {(0,1)(10,.64)(20,.43)(30,.32)(40,.27)(50,.25)};
        \addlegendentry{示例设置 A}
        \addplot[Teal,thick,dashed,mark=square*,mark size=1.5pt]
          coordinates {(0,1)(10,.51)(20,.32)(30,.24)(40,.21)(50,.20)};
        \addlegendentry{示例设置 B}
      \end{axis}
      \end{tikzpicture}
    \end{column}
    \begin{column}{.28\textwidth}
      \labeltext{图中可见}
      \vspace{2mm}
      \small
      两条示意曲线随迭代下降，后期变化趋缓。
      \vspace{5mm}

      真实报告需说明重复次数与不确定性。
    \end{column}
  \end{columns}
  \vspace{3mm}
  \smallnote{合成数据，仅展示图表版式。两条曲线不代表任何真实算法的性能。}
\end{frame}

% 06 / 表格：booktabs 三线表
\begin{frame}{对比表示例：表示维度与计算开销}
  \vspace{3mm}
  \begin{center}
    \renewcommand{\arraystretch}{1.55}
    \setlength{\tabcolsep}{7mm}
    \begin{tabular}{@{}lrrr@{}}
      \toprule
      \textbf{设置} & \textbf{维度 $d$} &
      \textbf{重构误差} & \textbf{耗时 / s}\\
      \midrule
      设置 A & 8  & 0.38 & 12.4\\
      设置 B & 16 & 0.29 & 18.7\\
      设置 C & 32 & 0.24 & 31.6\\
      \bottomrule
    \end{tabular}
  \end{center}
  \vspace{5mm}
  \labeltext{示例解读}
  在这组合成数值中，维度越高，误差越低，耗时越长。
  \vspace{5mm}
  \smallnote{合成数据，仅展示三线表。误差为无量纲示意值，耗时不来自实际测量。}
\end{frame}

% 07 / 总结与讨论
\section{结论与讨论}
\begin{frame}{结论与讨论}
  \vspace{2mm}
  \labeltext{模型表达}
  \par
  同一目标函数可以组合表达重构项与空间正则项。
  \vspace{5mm}

  \labeltext{参数权衡}
  \par
  $\lambda$ 控制两项的相对权重，$d$ 决定表示维度。
  \vspace{5mm}

  \labeltext{待验证的问题}
  \par
  空间图的构建方式会怎样影响结果？如何设计独立验证？
  \vfill
  \smallnote{以上为教学模型的讨论提纲，真实结论需要实际数据与实验支持。}
\end{frame}

% 08 / 参考文献：占位字段有意保留，供实际报告替换
\begin{frame}{参考文献}
  \smallnote{以下条目为格式占位，不是实际文献。正式报告时替换为真实来源。}
  \vspace{6mm}
  \begin{enumerate}
    \setlength{\itemsep}{1.2em}
    \item 作者 A，作者 B.\\
      \textit{论文标题}. 期刊或会议名称，年份.
    \item 作者 C，作者 D.\\
      \textit{论文标题}. 期刊或会议名称，年份.
    \item 项目或数据集名称.\\
      版本号，年份，访问地址.
  \end{enumerate}
\end{frame}

% 09 / 结束页
\begin{frame}[plain]
  \vfill
  {\color{Teal}\small DISCUSSION\par}
  \vspace{6mm}
  {\color{Navy}\LARGE\bfseries 谢谢，欢迎讨论\par}
  \vspace{5mm}
  {\color{Muted}\normalsize 汇报人姓名\quad 学校 / 课题组\par}
  \vspace{6mm}
  {\color{Teal}\rule{20mm}{1.2pt}\par}
  \vfill
\end{frame}

\end{document}
